\chapter{函数重载、命名空间和头文件边界}

\section{问题入口：同名函数为什么有时可以，有时不行}

\cpp 允许函数重载：同一个名字可以有不同参数列表。

\begin{lstlisting}[language=C++,caption={函数重载}]
int area(int side)
{
    return side * side;
}

int area(int width, int height)
{
    return width * height;
}
\end{lstlisting}

调用 \code{area(5)} 时选择第一个，调用 \code{area(5, 3)} 时选择第二个。重载的价值是让同一个概念共享名字。但如果两个函数语义不同，就不要硬用同名。名字相同应该意味着读者可以用同一种心智模型理解它们。

\section{重载解析从参数类型开始}

看这个例子：

\begin{lstlisting}[language=C++,caption={重载可能产生歧义}]
void print(long value);
void print(double value);

print(1);
\end{lstlisting}

\code{1} 是 \code{int}，既可以转换成 \code{long}，也可以转换成 \code{double}。编译器要根据转换等级选择候选；有时会报歧义。初学阶段不要设计过于相近的重载，尤其是整数、浮点、布尔和指针混在一起时。

\section{默认参数和重载不要打架}

下面的接口看起来方便，实际可能制造歧义：

\begin{lstlisting}[language=C++,caption={默认参数和重载混用风险}]
void connect(std::string_view host, int port = 80);
void connect(std::string_view host);
\end{lstlisting}

调用 \code{connect("example.com")} 时，两者都像是可行候选。更好的做法是只保留一个清楚接口，或者用配置对象表达可选项：

\begin{lstlisting}[language=C++,caption={用选项对象表达默认值}]
struct ConnectOptions {
    std::string host;
    int port = 80;
    bool use_tls = false;
};

void connect(const ConnectOptions& options);
\end{lstlisting}

当参数越来越多时，选项对象比一串默认参数更可读，也更容易扩展。

\section{命名空间控制名字范围}

大型项目不能把所有名字放在全局命名空间。命名空间给名字加上归属：

\begin{lstlisting}[language=C++,caption={命名空间组织模块}]
namespace score {

struct Summary {
    std::size_t count = 0;
    int sum = 0;
};

Summary summarize(std::span<const int> values);

} // namespace score
\end{lstlisting}

调用时写 \code{score::summarize(values)}，读者知道函数来自哪个模块。不要在头文件里写 \code{using namespace std;}，它会污染所有包含这个头文件的翻译单元，制造名字冲突。

\section{头文件放声明，源文件放实现}

普通非模板函数通常在头文件声明，在源文件实现：

\begin{lstlisting}[language=C++,caption={summary.hpp}]
#pragma once

#include <span>

namespace score {

struct Summary {
    std::size_t count = 0;
    int sum = 0;
};

Summary summarize(std::span<const int> values);

} // namespace score
\end{lstlisting}

\begin{lstlisting}[language=C++,caption={summary.cpp}]
#include <score/summary.hpp>

namespace score {

Summary summarize(std::span<const int> values)
{
    Summary result;
    result.count = values.size();
    for (int value : values) {
        result.sum += value;
    }
    return result;
}

} // namespace score
\end{lstlisting}

头文件是依赖传播边界。头文件包含越多，编译依赖越重。能在源文件包含的实现细节，就不要放进公开头文件。

\section{include guard 和 pragma once}

头文件可能被多次包含，需要防止重复定义。现代项目常用 \code{#pragma once}：

\begin{lstlisting}[language=C++,caption={防止头文件重复包含}]
#pragma once

// declarations
\end{lstlisting}

传统写法是 include guard：

\begin{lstlisting}[language=C++,caption={传统 include guard}]
#ifndef SCORE_SUMMARY_HPP
#define SCORE_SUMMARY_HPP

// declarations

#endif
\end{lstlisting}

两者目的相同。团队已有风格优先跟随。关键不是选哪个，而是所有头文件都要有保护。

\section{模板为什么常放头文件}

模板在使用点实例化。编译器看到 \code{average<int>} 的调用时，需要模板定义本身，而不只是声明。因此模板通常放在头文件：

\begin{lstlisting}[language=C++,caption={模板定义放头文件}]
template <typename T>
T max_value(T lhs, T rhs)
{
    return lhs < rhs ? rhs : lhs;
}
\end{lstlisting}

这也是模板代码会增加编译依赖的原因。模板越大，包含它的文件越多，编译时间越容易上升。泛型代码要更克制。

\section{完整案例：拆一个小数学库}

目录：

\begin{lstlisting}[caption={小库目录}]
mathlib/
  include/mathlib/statistics.hpp
  src/statistics.cpp
  tests/statistics_tests.cpp
\end{lstlisting}

公开接口只暴露必要类型和函数：

\begin{lstlisting}[language=C++,caption={statistics.hpp}]
#pragma once

#include <optional>
#include <span>

namespace mathlib {

struct AverageResult {
    double value = 0.0;
};

std::optional<AverageResult> average(std::span<const int> values);

} // namespace mathlib
\end{lstlisting}

测试只包含公开头文件，不包含 \code{src/statistics.cpp}。如果测试必须包含源文件才能访问内部函数，说明边界可能还没设计好。

\begin{exercisebox}
把成绩统计函数放进 \code{score} 命名空间，并拆成头文件、源文件、测试文件。要求头文件不包含 \code{iostream}，测试通过公开接口调用，不直接包含 \code{.cpp} 文件。
\end{exercisebox}
